% Options for packages loaded elsewhere
\PassOptionsToPackage{unicode}{hyperref}
\PassOptionsToPackage{hyphens}{url}
\PassOptionsToPackage{dvipsnames,svgnames,x11names}{xcolor}
%
\documentclass[
  letterpaper,
  DIV=11,
  numbers=noendperiod]{scrartcl}

\usepackage{amsmath,amssymb}
\usepackage{iftex}
\ifPDFTeX
  \usepackage[T1]{fontenc}
  \usepackage[utf8]{inputenc}
  \usepackage{textcomp} % provide euro and other symbols
\else % if luatex or xetex
  \usepackage{unicode-math}
  \defaultfontfeatures{Scale=MatchLowercase}
  \defaultfontfeatures[\rmfamily]{Ligatures=TeX,Scale=1}
\fi
\usepackage{lmodern}
\ifPDFTeX\else  
    % xetex/luatex font selection
\fi
% Use upquote if available, for straight quotes in verbatim environments
\IfFileExists{upquote.sty}{\usepackage{upquote}}{}
\IfFileExists{microtype.sty}{% use microtype if available
  \usepackage[]{microtype}
  \UseMicrotypeSet[protrusion]{basicmath} % disable protrusion for tt fonts
}{}
\makeatletter
\@ifundefined{KOMAClassName}{% if non-KOMA class
  \IfFileExists{parskip.sty}{%
    \usepackage{parskip}
  }{% else
    \setlength{\parindent}{0pt}
    \setlength{\parskip}{6pt plus 2pt minus 1pt}}
}{% if KOMA class
  \KOMAoptions{parskip=half}}
\makeatother
\usepackage{xcolor}
\setlength{\emergencystretch}{3em} % prevent overfull lines
\setcounter{secnumdepth}{-\maxdimen} % remove section numbering
% Make \paragraph and \subparagraph free-standing
\makeatletter
\ifx\paragraph\undefined\else
  \let\oldparagraph\paragraph
  \renewcommand{\paragraph}{
    \@ifstar
      \xxxParagraphStar
      \xxxParagraphNoStar
  }
  \newcommand{\xxxParagraphStar}[1]{\oldparagraph*{#1}\mbox{}}
  \newcommand{\xxxParagraphNoStar}[1]{\oldparagraph{#1}\mbox{}}
\fi
\ifx\subparagraph\undefined\else
  \let\oldsubparagraph\subparagraph
  \renewcommand{\subparagraph}{
    \@ifstar
      \xxxSubParagraphStar
      \xxxSubParagraphNoStar
  }
  \newcommand{\xxxSubParagraphStar}[1]{\oldsubparagraph*{#1}\mbox{}}
  \newcommand{\xxxSubParagraphNoStar}[1]{\oldsubparagraph{#1}\mbox{}}
\fi
\makeatother


\providecommand{\tightlist}{%
  \setlength{\itemsep}{0pt}\setlength{\parskip}{0pt}}\usepackage{longtable,booktabs,array}
\usepackage{calc} % for calculating minipage widths
% Correct order of tables after \paragraph or \subparagraph
\usepackage{etoolbox}
\makeatletter
\patchcmd\longtable{\par}{\if@noskipsec\mbox{}\fi\par}{}{}
\makeatother
% Allow footnotes in longtable head/foot
\IfFileExists{footnotehyper.sty}{\usepackage{footnotehyper}}{\usepackage{footnote}}
\makesavenoteenv{longtable}
\usepackage{graphicx}
\makeatletter
\def\maxwidth{\ifdim\Gin@nat@width>\linewidth\linewidth\else\Gin@nat@width\fi}
\def\maxheight{\ifdim\Gin@nat@height>\textheight\textheight\else\Gin@nat@height\fi}
\makeatother
% Scale images if necessary, so that they will not overflow the page
% margins by default, and it is still possible to overwrite the defaults
% using explicit options in \includegraphics[width, height, ...]{}
\setkeys{Gin}{width=\maxwidth,height=\maxheight,keepaspectratio}
% Set default figure placement to htbp
\makeatletter
\def\fps@figure{htbp}
\makeatother

\KOMAoption{captions}{tableheading}
\makeatletter
\@ifpackageloaded{caption}{}{\usepackage{caption}}
\AtBeginDocument{%
\ifdefined\contentsname
  \renewcommand*\contentsname{Table of contents}
\else
  \newcommand\contentsname{Table of contents}
\fi
\ifdefined\listfigurename
  \renewcommand*\listfigurename{List of Figures}
\else
  \newcommand\listfigurename{List of Figures}
\fi
\ifdefined\listtablename
  \renewcommand*\listtablename{List of Tables}
\else
  \newcommand\listtablename{List of Tables}
\fi
\ifdefined\figurename
  \renewcommand*\figurename{Figure}
\else
  \newcommand\figurename{Figure}
\fi
\ifdefined\tablename
  \renewcommand*\tablename{Table}
\else
  \newcommand\tablename{Table}
\fi
}
\@ifpackageloaded{float}{}{\usepackage{float}}
\floatstyle{ruled}
\@ifundefined{c@chapter}{\newfloat{codelisting}{h}{lop}}{\newfloat{codelisting}{h}{lop}[chapter]}
\floatname{codelisting}{Listing}
\newcommand*\listoflistings{\listof{codelisting}{List of Listings}}
\makeatother
\makeatletter
\makeatother
\makeatletter
\@ifpackageloaded{caption}{}{\usepackage{caption}}
\@ifpackageloaded{subcaption}{}{\usepackage{subcaption}}
\makeatother

\ifLuaTeX
  \usepackage{selnolig}  % disable illegal ligatures
\fi
\usepackage{bookmark}

\IfFileExists{xurl.sty}{\usepackage{xurl}}{} % add URL line breaks if available
\urlstyle{same} % disable monospaced font for URLs
\hypersetup{
  colorlinks=true,
  linkcolor={blue},
  filecolor={Maroon},
  citecolor={Blue},
  urlcolor={Blue},
  pdfcreator={LaTeX via pandoc}}


\author{}
\date{}

\begin{document}


\section{Main Thesis of GenBrain}\label{main-thesis-of-genbrain}

Cardiology made major strides when an analogy was drawn between the
heart and a mechanical pump. Descartes, Harvey, Hales, and others began
to understand the circulatory system by watching how a bellows moves air
or how piston-based hydraulic systems like fountains or firefighting
systems moved water from place to place. This critical analogy
eventually allowed for diagnoses using pump terminology (``valves'',
``ejection fraction'', ``blood pressure''), and predictions about heart
outcomes given analagous insults to a mechanical device (i.e.~making a
hole in a pump tube will induce the same outcome as a vascular
perforation; clogging the tubes of the pump requires more force, as in
hypertension). Nearly every aspect of cardiology is now framed in this
view, and as any good model should, the pump model opened doors for
framing questions and experiments, understanding the heart's failure
modes, and designing effective interventions and treatments.

Unfortunately, there is no useful analogy that has allowed us to develop
a similar multi-level understanding of the brain. We currently describe
mental health phenomena at the level of behavioral and thought patterns;
these patterns can be observed and categorized in order to prescribe
treatments that generally \emph{stabilize} affected individuals, but are
largely ineffective at halting recurring thought disorders. The organ
that generates these thinking patterns remains a mystery: comparison
between the brains of individuals with schizophrenia and those without
provides little insight, because brain ``abnormalities'' that are
associated with the disorder (e.g.~ventricle size, frontal cortex size,
hippocampus size) fall within the range of normal human variance.
Genetic abnormalities have likewise come up empty, again with most
genetic anomalies observed consistently in the healthy human population.
Without a reference point to gauge how brains actually work, and without
a clear smoking gun from non-computational hypotheses like anatomy and
genetics, we should turn to cardiology's example and answer ``What is
the pump of the brain?'' This is a difficult question because the
analogy is likely 3-pronged (see below), and our current best analogies
fail to provide inroads to understanding. For example, the ``digital
computer'' analogy fails because, unlike computers:

\begin{enumerate}
\def\labelenumi{\arabic{enumi}.}
\tightlist
\item
  The hardware of the brain is not fixed. It is constantly rewiring
  itself and changing throughout development to adapt to new
  environments and behavioral demands
\item
  The brain operates on about 10 watts of power and individual units
  operate at a max of \textasciitilde500Hz: GPU-based computers consume
  about two orders of magnitude more power, with units operating in the
  GHz range.
\item
  The brain is aware of itself and other agents. It undergoes subjective
  experience, moods, and emotions.
\item
  The brain learns mostly in an undirected environment without clear
  rules, and generalizes through analogy.
\end{enumerate}

Recently, deep neural networks have become a popular analogy for brains
due to their adaptibility, their use of neural-like elements, and their
extraordinary performance on human-inspired tasks including language and
vision. Some of our collaborators have taken the analogy quite
literally, proposing that layers of a deep ANN are akin to steps in the
ventral visual stream of the neocortex. However, ANNs are a poor analogy
for the workings of the brain, although they may be very useful tools
for AI. ANNs differ from real brains along a number of axes that
preclude a useful analogy like the heart-pump relationship.

\begin{enumerate}
\def\labelenumi{\arabic{enumi}.}
\item
  The brain is composed of varying neurotransmitter systems,
  neuromodulators, and neuron types. ANNs have only two ways of
  impacting downstream neurons - positive or negative weights.
\item
  Neurons and synapses in the brain are exceedingly stochastic. Activity
  for the same stimulus can vary wildly from trial to trial, while
  percepts and actions remain intact. ANNs perform deterministic
  transforms during inference, although some training methods inject
  randomness to improve generalizability during inference.
\item
  The brain is a highly recurrent system; feedforward models, and even
  RNNs with outputs recurring to inputs, do not capture the clear
  recursive patterns that exist in brain connectivity.
\item
  ANNs are a feedforward model that provides categorical outputs when
  fed data. But the brain's models are \emph{generative}, meaning we can
  generate data and lower level features conditioned on knowledge of
  high level constructs. ANNs would require a second artifact for this
  computation.
\end{enumerate}

\section{SMC with Inference Control}\label{smc-with-inference-control}

Instead of a broader ``computer'' analogy, GenBrain subscribes to the
idea that the brain is a special kind of computer running a particular
type of program that is capable of tracking latent states as they evolve
in real time, recognizing the uncertainty associated with state
inference, and using this uncertainty to guide adaptive choices and
learning. The paradigm is called SMC with Inference Control.

\begin{figure}[H]

{\centering \includegraphics{AllenBrain_SMCNN_w_Baseline.png}

}

\caption{SMC in the Brain}

\end{figure}%

SMC models and inference programs are written in specialized computer
languages; our group's specialty is one of these languages (GenJAX), and
we have made strides comparing the features and control flow of SMC with
real brain circuits. Below is our current conception about how core
features of cognition map to SMC models that we have been working on in
CHI, with each mapping have a neurobiological correlate that we can
attempt to explain with a unified model of cognition, computataion, and
neuroscience.

\begin{longtable}[]{@{}
  >{\centering\arraybackslash}p{(\columnwidth - 4\tabcolsep) * \real{0.4079}}
  >{\centering\arraybackslash}p{(\columnwidth - 4\tabcolsep) * \real{0.2763}}
  >{\centering\arraybackslash}p{(\columnwidth - 4\tabcolsep) * \real{0.3158}}@{}}
\toprule\noalign{}
\begin{minipage}[b]{\linewidth}\centering
Core Cognition
\end{minipage} & \begin{minipage}[b]{\linewidth}\centering
SMC + Inference Control
\end{minipage} & \begin{minipage}[b]{\linewidth}\centering
Neural Correlate
\end{minipage} \\
\midrule\noalign{}
\endhead
\bottomrule\noalign{}
\endlastfoot
Ability to survive and interact in the physical world & SMC object
perception, intuitive physics, MC action planning & Ventral stream,
MBVS \\
Ability to predict other people's thoughts and actions & SMC inverse
planning & Mirror neurons \\
Ability to localize experience to self in place and time & SLAM,
Self-Orienting Models & Place cells, grid cells \\
Ability to weigh actions based on value of outcome & Utility-based
particle reweighting & Dopamine system \\
Ability to learn new generative models and adjust inborn models & Open
Universe models, VI / Gradients, Boltzmann Learning & LTP / LTD,
Acetylcholine, Axon rerouting \\
\end{longtable}

Its important to note that each facet of this table is an active area of
research. Object perception, intuitive physics, SLAM, and planning are
all extremely difficult problems where state of the art AI systems
cannot reach the performance of the brain in any sense (performance,
efficiency, speed, flexibility, ease of learning). Understanding how the
brain achieves such incredible performance at these tasks should shed
light on AI research (e.g.~the feature space of likelihood, massive
parallelism, many largely fixed modules). Moreover, GenJAX is a perfect
language for pursuing these goals. GenJAX allows massively parallel
particle filtering and we are actively developing inference control
strategies and fixed depth recursive inference required for GenMind
modeling.

It is also important to note that each facet of cognition described here
has been shown to be dysregulated in particular mental health disorders.
A mechanistic neurocomputational description of neural dysfunction is
completely missing in neuropsychiatry, and its something Dost Ongur has
repeatedly said is up to us to solve so he can set actionable goals.

NEED TO ADD NEURAL COLLABORATIONS OUTSIDE OF JANELIA (wickersham,
wiktor; caroline runyan) Put specific people into the long term goals!
Put these things in the talk!

\section{Roadmap}\label{roadmap}

\subsubsection{Q4 2024 Effort}\label{q4-2024-effort}

\begin{enumerate}
\def\labelenumi{\arabic{enumi}.}
\item
  I will be meeting with Vikash, Josh, and Jim, starting on November
  18th, to outline a talk that can be delivered to neuroscientists with
  little background in computation, and to write an RDD describing the
  outline of a SMCNN manuscript.

  \begin{itemize}
  \tightlist
  \item
    \emph{Deliverables}: RDD for slides and paper
  \item
    \emph{Contributor Requirements}: \textasciitilde3-4 1 hour meetings
    with Jim, Josh, and Vikash. 1-2 hrs effort by ADB in preparing for
    these meetings and 4-5 hrs per week writing the RDD outside of these
    meetings. 2 hr meeting every 2 weeks with GM.
  \item
    \emph{Approximate Dates} : Nov 18 2024-January 2025
  \end{itemize}
\item
  For the SMCNN paper, we have decided to port the 2D-\textgreater3D dot
  perception model from the original SMCNN draft (Julia) to GenJAX. We
  also decided to make a live demo where users can enter their own 2D
  trajectories.

  \begin{itemize}
  \tightlist
  \item
    \emph{Deliverables}: A model for the Marquee Figure for the SMCNN
    paper
  \item
    \emph{Contributor Requirements}: 40+ hours of ADB coding / debugging
    time. Workshopping with GM for stimuli that
  \item
    \emph{Approximate Dates} : Oct 18 2024-January 2025
  \end{itemize}
\item
  We have decided to try to map the topology of SMCNNs to real neural
  circuits in brains, selecting cortical regions for each latent in the
  2D-\textgreater3D model. This idea will result in a figure of a
  transverse section of the brain with variable transmission mapped onto
  real axonal pathways in the Allen Brain Atlas.

  \begin{itemize}
  \tightlist
  \item
    \emph{Deliverables}: A figure for the SMCNN paper illustrating the
    mapping of the theory onto the real brain
  \item
    \emph{Contributor Requirements}: Difficult to estimate. The mouse
    brain is much different than the more heavily characterized monkey
    visual system.
  \item
    \emph{Approximate Dates} : Nov 8 2024-January 2025
  \end{itemize}
\item
  After completing 1-3, we will submit a first paper to bioarxiv, and
  simultaneous first-pass journal review. This will require new
  explanatory figures that will be detailed in the Part 1 RDD.

  \begin{itemize}
  \tightlist
  \item
    \emph{Deliverables}: Manuscript with updated figures
  \item
    \emph{Contributor Requirements}: 15 hrs per week ADB, 5 hrs per week
    GM. 2-3 hours editing and feedback from JT, JD, VKM, CR.
  \item
    \emph{Approximate Dates} : Nov15-Dec31 2024
  \end{itemize}
\item
  Preparing a talk for neuroscientists (which will require new figures
  and movies)

  \begin{itemize}
  \tightlist
  \item
    \emph{Deliverables}: A slide deck geared towards neuroscientists,
    using figures / material generated in Pt. 2
  \item
    \emph{Contributor Requirements}: 10 hrs per week effort ADB, 2-3 1hr
    feedback sessions with JT, VKM. Delivered in FRO or Probcomp
    Research Meetings.
  \item
    \emph{Approximate Dates} : Nov30-Dec31 2024
  \item
    \emph{Action Items} : Schedule FRO talk with Allison.
  \end{itemize}
\item
  Andrew, Vikash, Josh, and Tim write research product vision statements
  for implementing GenMind, testing of GenBrain predictions using
  neuroscience collaborations, and mental health exploratory projects.

  \begin{itemize}
  \tightlist
  \item
    \emph{Deliverables}: 2-page descriptions for each vision statement,
    with software-level Linear card descriptions of tasks to be
    performed for each point.
  \item
    \emph{Contributor Requirements}: VKM, Josh, Tim \textasciitilde4-5 1
    hr meetings, 5 hrs per week ADB. ADB will track scheduling.
  \item
    \emph{Approximate Dates} : Initial brainstorming meeting by December
    15, where we make high level choices about impact towards basic
    vs.~translational neuroscience.
  \end{itemize}
\item
  Janelia planning and initial setup of virtual environment

  \begin{itemize}
  \tightlist
  \item
    \emph{Deliverables}: A virtual game where actions influence the
    control of an agent intercepting a ball in a field of barriers,
    ramps, and occluders.
  \item
    \emph{Contributor Requirements}: ADB 30min-1hr per week to keep
    collaboration going until 1-3 are complete.
  \end{itemize}
\item
  Theory hour and theory hour prep, 2025 planning + roadmapping,
  mentoring, initial development of new SMCNN infrastructure (GenSP
  focused - see Q1 2025).

  \begin{itemize}
  \tightlist
  \item
    \emph{Deliverables}: Genjax will have SMCNNs as a distribution type,
    and SMCNNs will become a component of GenJAX rather than a separate
    library.
  \item
    \emph{Contributor Requirements}: 8hr per week ADB, 4hr for MH.
  \end{itemize}
\end{enumerate}

\subsection{Q1 to Q2 2025 - Initial Execution of Product Roadmaps from
Q4 2024
(4).}\label{q1-to-q2-2025---initial-execution-of-product-roadmaps-from-q4-2024-4.}

-- here, want to add ``person i want to hire and what would they do?''

After the first phase of SMCNN exposure via the talk and paper in Q4, we
will begin to execute the plans designed in Pt. 4 above and integrate
that work into our collaboration with Rob Johnson and Janelia. These are
the specific projects that will receive roadmapping attention from Pt 4.
Our models will have to self-localize the rat agent within a complex 3D
environment, localize 3D objects that occlude ball paths and determine
goal locations, and predict physical trajectories that targets will take
through space. In a later stage (Q3) we will also plan for the goals of
other rat agents, learn new features of the environment online, and
understand reward contingencies as they evolve within the task.

\begin{enumerate}
\def\labelenumi{\arabic{enumi}.}
\item
  Build out GenJAX models that frame each core cognition aspect as SMC
  with Inference Control (\emph{GenMind})

  \begin{itemize}
  \item
    3D object perception in the RatFetch Environment (ramps, balls,
    occluders, collidable barriers). Influenced by real neural sturcture
    (from theory hour) with ability to make neural predictions (retina,
    tectum, VC).

    \begin{itemize}
    \tightlist
    \item
      \emph{Deliverables}: A retina-based model that allows 2D feature
      extraction and feature-based likelihood computation in a dynamic
      environment.
    \item
      \emph{Contributor Requirements}: ADB 20\%, RJ 30\%, Theory Hour
      Attendees (5\%).
    \item
      \emph{Applications to GenBrain Goals}: First attempt at informing
      AI with neuroscience ideas
    \end{itemize}
  \item
    Intuitive physical reasoning. Build out infrastructure in Arijit's
    2D Red Green task for eventual applicability to 3D worlds with same
    contingencies.

    \begin{itemize}
    \tightlist
    \item
      \emph{Deliverables} : A first paper with traceable physics in
      GenJAX, illustrating our approach to building out a system for
      rejuvenation, resampling, reasoning within a 2D world.
    \item
      \emph{Contributor Requirements} : AD 75\%, KS 20\%, AB 20\%
    \end{itemize}
  \item
    2D and 3D SLAM for self-localization and object localization within
    the Janelia RatFetch environment

    \begin{itemize}
    \tightlist
    \item
      \emph{Deliverables} : An SMCNN compilable SLAM model of
      localization within the Janelia environment, using first person
      observations as sensor data, and robot-driven novel barrier
      contingencies to show mapping rather than learned paths.
    \item
      \emph{Contributor Requirements} 20\% ADB, 5\% MH, 5\% GM. MK?
    \end{itemize}
  \end{itemize}
\item
  Continued progress on GenSP use in SMCNN integration to GenJAX

  \begin{itemize}
  \tightlist
  \item
    \emph{Deliverables} : Ability to apply SMCNN easily to any GenJAX
    model.
  \item
    \emph{Contributor Requirements} : AB 20\%, MH 20\%, MB 5\%
  \end{itemize}
\item
  Continued mentorship / theory hour development, roadmapping for future
  quarters, visitation to Janelia for continued collaboration.

  \begin{itemize}
  \tightlist
  \item
    \emph{Contributor Requirements} : AB 20\%, RJ 20\%, MH 10\%.
  \end{itemize}
\end{enumerate}

\subsection{Q3 2025+Later.}\label{q3-2025later.}

\begin{enumerate}
\def\labelenumi{\arabic{enumi}.}
\tightlist
\item
  Build inference controllers to influence inference patterns in Pt 1.
  models.

  \begin{itemize}
  \tightlist
  \item
    Inference control examples:

    \begin{itemize}
    \tightlist
    \item
      2D mapped likelihood inducing attentional shifts to unexplained
      regions of space
    \item
      Generate more particles when likelihood scores are low across
      current particle set
    \item
      Stop inference, replay and rejuvenate past hypothesis, continue
      inference.
    \item
      Realize model is failing, and that no other models are working.
      Learn about the world.
    \item
      Update reward contingencies when task variables change.
    \end{itemize}
  \end{itemize}
\item
  Buildout of SMCNN circuits to accomodate development of all GenMind
  features above, using GenSP based SMCNN computations within GenJAX
  that use biological neurons as untraced random choices.

  \begin{itemize}
  \tightlist
  \item
    Implement SMCNNs in more biologically realistic neurons and examine
    limitations of non-ideal neurons. Possible use of Nengo, but open to
    any suggestions.
  \item
    Generate novel roles for neurotransmitter systems in inference
    control.
  \end{itemize}
\item
  Use Janelia Rat Fetch task as a testbed for models of animal cognition

  \begin{itemize}
  \tightlist
  \item
    Planning in single rat experiments for single ball interception,
    Inverse planning in multi-agent experiments with other rats (who are
    also trying to get the ball)
  \item
    Learning of new object contingencies that influence value of rewards
    (i.e.~more milkshake for particular ball colors or sizes, new
    objects)
  \item
    Neural recordings in all hypothesized brain regions of SMCNNs.
  \end{itemize}
\end{enumerate}

\section{Related but tangential
projects}\label{related-but-tangential-projects}

\begin{itemize}
\item
  Jim Dicarlo Collaboration (``Whats Behind Task''). Jim remains very
  interested in SMCNNs and testing them. I don't have a lot of
  confidence that a monkey-based research program will yield insight,
  but I'm excited to work with Jim as a collaborator generally on
  psychophysics and neuroscience ideas.
\item
  GenBrain Score GenBrain score is a paradigm similar to brain score but
  extending it to accomodate spiking models, dynamic models, and
  uncertainty representation. Wiktor Phillips is in charge of this
  project and he is kind of on ``probation'' to see if he can
  independently drive research projets forward, or is more of an
  engineer.
\item
  GenBrain Atlas Pierre Glaser worked with us on modeling how SMCNNs
  could be used to predict model structure in the Allen Brain Atlas.
  This project is conceptually vague and Pierre's modeling strategy is
  new to me.
\end{itemize}

\section{\texorpdfstring{Applications to mental health, end 2025
(\emph{The EASE scale of anomalous self
experience})}{Applications to mental health, end 2025 (The EASE scale of anomalous self experience)}}\label{applications-to-mental-health-end-2025-the-ease-scale-of-anomalous-self-experience}

THe EASE scale is a rating of Anomalous Self Experience. High scores on
the EASE scale have proven to specifically differentiate thought
disorders like schizophrenia from other forms of major mental illness
(i.e.~bipolar, major depression). The EASE scale incorporates multiple
diagnoses that I believe SMCNNs are in a position to explain through the
lens of probabilistic programming with SMC, SMCNNs, and reasoning under
uncertainty.

1.9 \emph{Ambivalence} (A.5) can be framed as the inability to properly
weight particles against each other, leading to a uniform particle
distribution where no hypothesis feels ``more correct'' than any other.
Interestingly, this is a prediction of dopamine dysregulation in SMCNNs
-- dopamine is the primary neurotransmitter associated with
schizophrenia in neurobiology.

Summary: Inability to decide between two or more options. Persistent and
painful conscious coexistence of contradic- tory inclinations or
feelings. Ambivalence occurs even for very simple or trivial everyday
decisions. The patient cannot decide at all, needs more time for his
decision, or becomes immediately uncertain about a finally made de-
cision and changes it again. A related phenomenon that is scored here is
when the patient complains of having contradictory thoughts or feelings
at exactly the same time.

Examples: • She has difficulty in making decisions because she
`considers things in many ways'. Yesterday, it took her 3 h to decide on
which gift to buy for her boyfriend. • At the teachers' college, she
reversed her choice of subjects three times but still couldn't make out
whether she had made the right choice. • He is `snowed under with
options'; e.g.~he thinks that he probably ought to become a vegetarian
even though he loves meat. Such considerations lead him into
`doubleness' and `silly, blind alleys'. • Each time I think of
something, I get a counterthought on the other side of the brain
{[}score here also spatialization of experience (1.8){]}.

1.1 \emph{Thought Interference} (C.1.1)

Summary: Contents of consciousness (thoughts, imaginations, or
impulses), semantically disconnected from the main line of thinking,
appear automatically (not necessarily quickly or many), break into the
main line of thinking and interfere with it. Such thoughts are often
(but not always) emotionally neutral and they do not need to have a
special or extraordinary meaning. The patient may use private
designations to describe such thoughts (`thought tics', `acute
thoughts', and `surrealistic thoughts').

There are a number of these types of thought dysregulations where it
feels like the answer might be ``simulating from the prior'' rather than
constraining on data. This could also be a case of uniform scoring,
where all hypothesis seem equally likely to explain the current data.

2.2 \emph{Distorted First-Person Perspective}

The sense of self as the absolute experiential point of orientation
{[}i.e.~as something which does not itself have a precise location (me
who is here; the self identical with all experiencing) but to which
everything else is spatially related (ego-centric space){]} or as a
pole/source/focus of experience or action (I-consciousness) may be felt
to be at a specific spatial location or to have characteristics of
extension, or sometimes being spatially dislocated {[}in both cases,
always rate also spatialization of experience (1.8){]}. (related to
I-split,

• I constantly regard myself. Sometimes it is so pronounced that I can
hardly follow what's going on on TV. Even during a con- versation with
others, I observe myself to the point of having difficulty in grasping
what my interlocutors are saying. • My own `I', as a point of
perspective, feels as if it had shifted a few centimeters backwards.

The GenBrain team thinks that there are multiple avenues for pursuing
this phenotype. SLAM models where the individual is simultaneously
localizing themselves and mapping themselves in the world feel like an
appropriate description of ``regarding oneself''. Tiffany is very
interested in this avenue, hoping to mix in her expertise in VR. Our
SLAM models that we make for rats will make firm predictions about place
representation, and place representation is the most ``believable''
result from neuroscience (i.e.~place, grid cells).

Interestingly, this ``fragmentation'' of self seems to carry over to
visual perception, where scenes are also fragmented into individual
components rather than unified wholes.

``She remembered that she could not look at the whole door. She could
only see the knob or some corner of the door. The wall was fragmented
into parts'' (ARIETI 1962, p.85);

``I may look at the garden, but I don't see it as I normally do. I can
only concentrate on detail. For instance, I can lose myself in looking
at a bird on a branch, but then I don't see anything else'' (MATUSSEK
1987, p.92);

``Everything I see is split up. It's like a photograph that's torn in
bits and put together again. If somebody moves or speaks, everything I
see disappears quickly and I have to put it together'' (CHAPMAN 1966,
p.29).

\subsection{Notes from Roadmapping
Meetings:}\label{notes-from-roadmapping-meetings}

\begin{enumerate}
\def\labelenumi{\arabic{enumi}.}
\item
  Is our vision to make a neural model that guides basic neuroscience
  research, or are we signposting treatments to translational
  neuroscienctists. Which one of these we pick is going to influence our
  most important decisions. Justify the commonality -- if its justified,
  you can expand the team b/c there will be overlap.
\item
  Nobody that uses deep ANNs cares about neuroscience, but the model
  emerged from it. Are SMCNNs anyway relevant in the same way? Yes!
  Because the brain is actually better at visual cognition than any AI
  systems.
\end{enumerate}




\end{document}
